\documentclass[conference]{IEEEtran}
\IEEEoverridecommandlockouts
% Template version as of 6/27/2024
% ---------------------------------------------------------------------------
% LaserLearn -- course paper.
%
% Scope: this paper reports a real-time laser-projection game platform that was
% built and field-tested, the two games we implemented on top of it (a hidden-
% thief hunt and a laser-slicing game), and the deployment work that made it run
% on a GPU-less laptop. The thief's adversary is presented in two forms that we
% actually built and compare: a trained tabular Q-learning policy and an A*
% pathfinding controller. Performance figures for the vision pipeline are taken
% from run logs and offline validation on the demonstration hardware (Sec. VII).
% ---------------------------------------------------------------------------

\usepackage{cite}
\usepackage{amsmath,amssymb,amsfonts}
\usepackage{algorithmic}
\usepackage{graphicx}
\usepackage{textcomp}
\usepackage{xcolor}
\usepackage{array}
\def\BibTeX{{\rm B\kern-.05em{\sc i\kern-.025em b}\kern-.08em
    T\kern-.1667em\lower.7ex\hbox{E}\kern-.125emX}}
\begin{document}

\title{LaserLearn: A Real-Time Green-Laser Vision Platform for Projected Games
on Commodity Hardware}

% --- AUTHOR BLOCK: explicit 2 / 2 / 1 grid, last author centered on row 3 ---
% Department label: adjust if your program is listed differently.
\newcommand{\authblk}[3]{%
  \begin{minipage}[t]{0.42\textwidth}\centering
  {\normalsize #1}\\[2pt]
  {\footnotesize
   \textit{Dept. of Computer Science and Engineering}\\
   \textit{United International University}\\
   ID: #2, Dhaka, Bangladesh\\
   #3\par}
  \end{minipage}%
}
\author{%
\begin{tabular}{@{}c@{\hspace{2.5em}}c@{}}
\authblk{Tonmoy Bishwas}{0152410030}{tbishwas2410030@bsds.uiu.ac.bd} &
\authblk{Samiha Mahzabin}{0152410071}{smahzabin2410071@bsds.uiu.ac.bd} \\[2em]
\authblk{Naimur Rahman}{0152410029}{nrahman2410029@bsds.uiu.ac.bd} &
\authblk{Asif Hasan Shanto}{0152410035}{ashanto2410035@bsds.uiu.ac.bd} \\[2em]
\multicolumn{2}{c}{\authblk{Lailatul Noor}{0152410041}{lnoor2410041@bsds.uiu.ac.bd}}
\end{tabular}%
}

\maketitle

\begin{abstract}
We present LaserLearn, a real-time platform that turns a handheld green laser
pointer into a game controller for a projected display, and the two games we
built on top of it. A phone camera films the projection; a single-class YOLO
detector, fine-tuned on our own green-laser dataset, finds the laser dot in each
frame (test mAP50 of 0.88 at full resolution); and a one-time four-corner
homography maps the dot into game space. Everything below the detector is
deliberately small: games receive the laser through one narrow input object that
exposes only what a real laser can do, which is what let us build a second,
quite different game without touching the vision code. The first game hides a
thief in the dark and uses the laser as a flashlight; the second is a
laser-slicing arcade game. For the thief's behaviour we built and compare two
adversaries: a tabular Q-learning policy trained against scripted laser
opponents (95 percent win rate in simulation), and an A* pathfinding controller
that plays on procedurally generated mazes. The harder problem was not accuracy
but speed on the demonstration machine, a laptop with no discrete GPU. Our main
practical finding is that on that machine real-time feel was limited by
contention for the CPU between the inference thread and the rendering loop rather
than by raw inference latency; moving inference onto the integrated GPU through
OpenVINO removed the contention and gave the largest single improvement we
measured. We close with the further games the same platform is meant to host.
\end{abstract}

\begin{IEEEkeywords}
laser-pointer interaction, object detection, YOLO, real-time inference, edge
inference, integrated-GPU offload, projector-camera systems, spatial augmented
reality, reinforcement learning, pathfinding
\end{IEEEkeywords}

% ===========================================================================
\section{Introduction}
% ===========================================================================
Most students meet computer vision and reinforcement learning through a screen.
A detector draws boxes on a recorded clip; an agent plays inside a simulator. The
techniques work, but they stay behind glass. We wanted to pull them out into a
single physical loop that a person can stand in front of and play, and to do it
with hardware a student actually owns.

LaserLearn is that loop. A projector throws the game onto a wall. A phone camera
on a tripod watches the wall and streams to a laptop. The player points an
ordinary green laser at the projection, and in every camera frame a fine-tuned
detector locates the dot and a homography carries its position into the game.
What the player does with that dot depends on the game. In our first game the
screen is almost entirely black and the laser is a flashlight: it reveals a small
circle of the world, and somewhere in the dark a thief is creeping through a maze
toward a treasure. Hold the light on the thief long enough and you catch it; let
it reach the treasure and it wins. Our second game uses the same laser as a
blade, and the player slices fruit that is thrown up across the screen.

Building two games this different on one pipeline was a deliberate goal, not an
afterthought. Underneath both sits the same chain---camera, detector, homography,
and a tiny input layer---and the games never see anything below that layer. They
receive one object each frame that reports where the laser is, whether it is
currently on the surface, and how fast it is moving. Nothing else. Because that
object is all a real laser can honestly provide, a game written against it cannot
accidentally depend on a mouse click or a key press, and a new game is mostly a
matter of deciding what to do with a moving point. The second game took far less
work than the first for exactly this reason.

This paper reports three things. First, the platform and its detection front
end---a green-laser detector we trained ourselves, the calibration that places
its output on the screen, and the input abstraction that the games sit on.
Second, the two games, including two different controllers we built for the
thief and what we learned by comparing them: a tabular Q-learning policy that
learns to evade in a fixed maze, and an A* controller that plans on a new random
maze every round. Third, and the part we found most interesting, the engineering
that made the whole thing feel instant on a GPU-less laptop.

That last point shaped the project more than any algorithm did. Our development
desktop has a discrete GPU and runs everything without complaint. The machine we
had to demonstrate on does not: it is a consumer laptop with an Intel processor
and integrated graphics. Closing that gap turned up a result we did not expect.
On the CPU-only target, the inference thread and the rendering loop were fighting
over the same processor, and the system felt sluggish even though each task on
its own was fast enough. Raw inference speed was a red herring. Moving inference
onto the integrated GPU through OpenVINO took it off the CPU entirely, ended the
contention, and was the biggest speed-up we got.

\noindent\textbf{Contributions.}
\begin{itemize}
\item A real-time, multithreaded laser-projection platform built around a
fine-tuned single-class green-laser detector and a four-corner homography, with a
narrow laser-faithful input layer that lets new games be added without touching
the vision pipeline.
\item Two games on that platform---a flashlight-reveal thief hunt and a
laser-slicing arcade game---and, for the thief, two adversary controllers we
built and compare: a trained tabular Q-learning policy and an A* planner.
\item A deployment study on commodity, GPU-less hardware whose central finding is
that CPU contention between inference and rendering, not inference speed, governs
real-time feel, and that integrated-GPU offload resolves it.
\end{itemize}

The rest of the paper is organised as follows. Section~\ref{sec:related} reviews
related work. Section~\ref{sec:system} describes the platform and its input
layer. Section~\ref{sec:methods} gives the methods: detection, calibration, the
flashlight reveal, and the two thief controllers. Section~\ref{sec:games}
describes the two games and the ones we plan next. Section~\ref{sec:challenges}
covers the implementation and deployment problems, Section~\ref{sec:results}
reports results, Section~\ref{sec:discussion} discusses limitations, and
Section~\ref{sec:conclusion} concludes.

% ===========================================================================
\section{Related Work}\label{sec:related}
% ===========================================================================
LaserLearn sits at the meeting point of several lines of work. We take each in
turn and end every part by saying how our system departs from it.
Table~\ref{tab:compare} summarises the comparison.

% ---------------------------------------------------------------------------
\subsection{Laser-Pointer Interaction and Laser-Dot Detection}
% ---------------------------------------------------------------------------
Tracking a laser pointer on a projected display with a camera is an old idea.
Kirstein and M\"uller showed a camera-tracked laser-pointer-as-mouse for
projection screens~\cite{kirstein1998}, and Olsen and Nielsen's ``Laser Pointer
Interaction''~\cite{olsen2001} is the standard reference: a camera looks for the
brightest red spot on the surface, a calibration step maps camera coordinates to
display coordinates, and the interaction copes with the laser's jitter and with
the lack of a real click. Later work pushed the idea to large and tiled
displays~\cite{ahlborn2005} and studied detection on its own. Lapointe and Godin
argued that reliable spot detection is the limiting factor for a usable
laser-pointer interface and built a fast classical detector~\cite{lapointe2005}.
For harder, less constrained scenes the detectors get more elaborate:
Krsti\'ni\'c et al. combine a modified circular Hough transform with optical-flow
motion analysis to pull the true dot out of competing bright blobs under a moving
camera~\cite{krstinic2014}, and Ch\'avez et al. mix template matching with a fuzzy
rule base to find a laser spot in real homes~\cite{chavez2012}. The shared
assumption, and the shared weakness, is that the dot is the brightest or reddest
thing in view---which fails under projector glare and coloured, moving content.

\textbf{How LaserLearn differs.} We drop the brightness-and-colour heuristic for a
single-class detector fine-tuned on the dot's learned appearance, which holds up
better against glare and against the animated content of a real game. Our camera
is a handheld phone, not a fixed calibrated rig, and the laser is not a cursor but
a spatial flashlight (and, in our second game, a blade)---uses this literature
does not cover.

% ---------------------------------------------------------------------------
\subsection{Projector-Camera Games and Spatial Augmented Reality}
% ---------------------------------------------------------------------------
Turning ordinary surfaces into interactive displays with a projector and a camera
is the premise of spatial augmented reality (SAR). Raskar et al. laid the
groundwork with the Office of the Future~\cite{raskar1998office}, the framing of
spatially augmented reality~\cite{raskar1998sar}, and Shader Lamps, which projects
image-based shading onto neutral objects~\cite{raskar2001shaderlamps}. Pinhanez's
Everywhere Displays steered a projector onto arbitrary surfaces and read hand
gestures with a camera, and was used to build projected games~\cite{pinhanez2001}.
More recent projector-camera gaming systems such as
IllumiRoom~\cite{jones2013illumiroom} and RoomAlive~\cite{jones2014roomalive}
extend play into a whole room using calibrated depth cameras and several projector
units. Closer to our budget, Song et al. built a tangible game interface from a
cheap projector-camera setup~\cite{song2007tangible}. The camera-to-surface map is
a homography; careful projector-camera calibration has been studied
thoroughly---Moreno and Taubin note that one global homography cannot model
projector lens distortion and use local homographies for sub-pixel
accuracy~\cite{moreno2012}, building on the standard treatment of planar
homographies~\cite{hartley2004}. The flashlight-as-reveal idea also has a
precedent: Ridel et al.'s Revealing Flashlight lets a user sweep a virtual light
over a heritage artifact to bring out surface detail~\cite{ridel2014}, itself a
descendant of the movable Magic Lens of Bier et al.~\cite{bier1993}.

\textbf{How LaserLearn differs.} We use the most minimal procam setup we
can---one consumer projector, one phone camera, one global four-corner
homography---rather than the multi-unit depth rigs of IllumiRoom and RoomAlive,
because our surface is flat and a coarse game grid does not need sub-pixel
registration. That same simplification, in line with the distortion caveat
of~\cite{moreno2012}, also explains a small spotlight offset we see and treat as a
known limitation (Section~\ref{sec:discussion}). Unlike the Revealing Flashlight,
our reveal is driven by a tracked physical laser rather than a tracked stylus, and
it is the core mechanic of a live game, not a tool for inspecting a static object.

% ---------------------------------------------------------------------------
\subsection{Real-Time Vision Deployment on Commodity Hardware}
% ---------------------------------------------------------------------------
Our detector is from the YOLO family of single-shot real-time
detectors~\cite{redmon2016}, whose later versions and small-object behaviour are
surveyed by Terven et al.~\cite{terven2023}; a laser dot is a very small object,
a case the original YOLO already flags as hard. Running such models off the GPU is
supported by toolkits like OpenVINO, whose analysis-and-tuning workbench targets
Intel CPUs and integrated GPUs~\cite{gorbachev2019}, and recent work benchmarks
inference frameworks on edge hardware~\cite{edgeframeworks2025}. A separate line of
systems work treats \emph{resource contention} as a first-class problem when
several models share one constrained device and schedules around
it~\cite{edgecontention2022}.

\textbf{How LaserLearn differs.} The benchmark literature mostly measures
inference \emph{in isolation}, and the contention literature looks at
\emph{multiple models} sharing a device. Our binding constraint is neither:
contention between a \emph{single} inference pipeline and an \emph{interactive
rendering loop} on a CPU-only consumer laptop. Our fix---pushing inference onto an
Intel Iris Xe integrated GPU via OpenVINO~\cite{gorbachev2019} to free the CPU for
rendering---turns that inference-versus-render trade-off from the dominant
bottleneck into a non-issue, which an isolated benchmark would never reveal.

% ---------------------------------------------------------------------------
\subsection{Adversary Behaviour: Learned and Search-Based}
% ---------------------------------------------------------------------------
We built the thief's behaviour two ways, so we draw on two literatures. For the
learned controller, Watkins and Dayan's Q-learning supplies the algorithm and a
convergence guarantee in the tabular, finite, discrete case~\cite{watkins1992},
and Sutton and Barto ground the MDP formulation, discounting, and exploration
schedule we use~\cite{sutton2018}. Deep RL---epitomised by the Deep Q-Network that
learns Atari from pixels~\cite{mnih2015}---is the tool of choice when the state is
large or continuous, which ours is not. Pursuit and evasion are naturally
multi-agent: Tan contrasts independent and cooperative Q-learners in a
predator-prey grid~\cite{tan1993}, Littman frames zero-sum play as Markov games
solved by minimax-Q~\cite{littman1994}, and more recent work learns pursuit
policies with deep RL and transfers them to real drones~\cite{desouza2021}. For
contrast, OpenAI's hide-and-seek shows competitive self-play producing an emergent
curriculum of strategies~\cite{baker2020}, and RL has also been used to make game
bots less predictable~\cite{glavin2015}. For the search-based controller we use
A*, the classic best-first graph search with an admissible heuristic of Hart,
Nilsson, and Raphael~\cite{hart1968}; in our grid the natural heuristic is
Manhattan distance, and the laser enters as an added traversal cost rather than as
a co-learning agent.

\textbf{How LaserLearn's adversaries differ.} Our state for the learned thief is a
compact four-tuple---thief cell and laser cell on a fifteen-by-nine grid---with
four actions, so a Q-table is small, trains in seconds on a CPU, is fully
readable, and inherits the convergence guarantee of~\cite{watkins1992}; deep RL as
in~\cite{mnih2015,desouza2021} would add instability and tuning pain for no gain at
this size. Rather than the self-play autocurriculum of~\cite{baker2020}, we train
against an \emph{exogenous}, hand-designed set of scripted laser behaviours,
following the independent-learner view of~\cite{tan1993}, which keeps training
cheap and reproducible and aimed at the opponents the game can actually present.
The A* controller takes a different stance entirely: it does no learning, replans
every cell, and so handles a freshly generated maze each round, where the trained
table is tied to the one maze it learned. Comparing the two on the same game is
part of what we report (Section~\ref{sec:adversaries}).

% ---------------------------------------------------------------------------
\subsection*{Research Gap}
% ---------------------------------------------------------------------------
Every ingredient has precedent: camera-tracked laser interaction is decades old,
procam games and SAR are mature, edge inference is active, and both tabular
Q-learning and A* are textbook. What we put together is the combination: a
\emph{fine-tuned single-class laser detector} feeding a \emph{homography-mapped
input layer} that carries one physical laser into \emph{more than one real-time
projected game} on \emph{commodity, GPU-less hardware}, with the deployment
finding that processor contention---not inference speed---sets the real-time feel.

\begin{table*}[t]
\caption{Representative Prior Systems vs. LaserLearn}
\label{tab:compare}
\begin{center}
\renewcommand{\arraystretch}{1.25}
\begin{tabular}{|p{2.6cm}|p{2.4cm}|p{2.7cm}|p{2.2cm}|p{2.2cm}|p{2.0cm}|}
\hline
\textbf{System} & \textbf{Input modality} & \textbf{Dot/target detection} &
\textbf{Calibration} & \textbf{Adversary} & \textbf{Hardware target} \\
\hline
Olsen \& Nielsen~\cite{olsen2001} & Laser pointer (cursor) & Brightest-red-spot
heuristic & Camera-to-display map & None & Desktop + camera \\
\hline
Pinhanez (Everywhere Displays)~\cite{pinhanez2001} & Hand gesture & Vision gesture
detection & Oblique-projection correction & None & Steerable projector rig \\
\hline
RoomAlive~\cite{jones2014roomalive} & Body / depth & Depth sensing & Self-calibrating
procam units & Scripted & Multiple projector-depth units \\
\hline
Revealing Flashlight~\cite{ridel2014} & Tracked stylus & Tracked device pose &
SAR projector calibration & None (inspection) & Workstation + projector \\
\hline
Hide-and-seek (OpenAI)~\cite{baker2020} & Simulated agents & N/A (simulator) & N/A &
Self-play (emergent) & Large-scale GPU cluster \\
\hline
\textbf{LaserLearn (ours)} & \textbf{Green laser pointer (flashlight + blade)} &
\textbf{Fine-tuned single-class YOLO} & \textbf{Single 4-corner homography} &
\textbf{Q-learning + A* thief} & \textbf{GPU-less laptop + iGPU} \\
\hline
\end{tabular}
\end{center}
\end{table*}

% ===========================================================================
\section{Platform and Input Layer}\label{sec:system}
% ===========================================================================
LaserLearn runs as a multithreaded pipeline whose stages are decoupled so that
slow inference never stalls rendering. A \emph{capture thread} reads only the
newest frame from the phone camera (camera buffer size one, no queue). A
\emph{detector thread} runs the YOLO model on that frame, keeps the
highest-confidence detection, and maps its box centre through the homography into
game coordinates, publishing the result as the current laser position. A
\emph{rendering loop} reads whatever the latest laser position is and draws the
frame. Because the detector only ever publishes the most recent position and the
renderer never waits on it, the detection rate and the render rate float free of
each other---which is exactly what makes the contention behaviour in
Section~\ref{sec:results} both visible and tunable. A device \emph{profile} fixes
the inference backend, image size, confidence threshold, and camera index before
the threads start, so the same code runs on the discrete-GPU desktop and the
integrated-GPU laptop with a one-word change.

\noindent\textbf{The input layer.} Sitting on top of all of this is one small
object, and it is the part of the design we lean on most. Every game receives, per
frame, a pointer snapshot: an $(x,y)$ position in the fixed
$1920\times1080$ game canvas, a boolean saying whether the laser is on the surface
this frame, and the velocity $(v_x,v_y)$ since the last frame (from which a speed
is derived). That is the entire contract. There is no mouse button, no key, no
scroll, no second cursor, because a laser pointer offers none of those. During
development we feed this same object from the mouse so the games can be built and
played on a normal computer; at the installation a laser-backed source fills it
instead, and the active flag simply follows whether the beam is currently being
detected. The games cannot tell the difference, and neither can their code.

Two helpers turn that bare pointer into usable interaction without ever adding a
click. A \emph{dwell tracker} fills a progress value from zero to one while the
pointer is held near a target and drains it (at half speed) when it drifts off;
crossing one fires the action. This is how the laser ``clicks''---to choose a game
from the menu, and, in the thief game, to make a catch. A \emph{gesture} helper
asks whether the straight segment between the pointer's previous and current
position passes within some radius of a target while moving above a speed
threshold; that is how the laser ``slices'' in the second game. Both helpers read
only the pointer snapshot, so any game can use them, and a new game that needs
``hover to select'' or ``swipe through'' gets it for free. We come back to why
this matters for adding games in Section~\ref{sec:moregames}.

% ===========================================================================
\section{Methods}\label{sec:methods}
% ===========================================================================

% ---------------------------------------------------------------------------
\subsection{Laser Detection}\label{sec:detect}
% ---------------------------------------------------------------------------
The laser dot is found by a single-class object detector. We use a green
532\,nm laser---bright enough to read on the projection under normal room
light---and fine-tuned a nano-scale YOLO model (about 2.5\,M parameters, the
smallest of a recent YOLO release) on our own dataset, annotated with one class,
\texttt{laser}. The dataset is our own: 1069 full-resolution
($1920\times1080$) frames, with 675 laser annotations across the single class. Two
choices in how we built it mattered. First, we shot the green dot against a wide
range of background colours rather than one neutral wall, so the detector would not
quietly learn ``green on grey'' and then fail when a game drew a coloured scene.
Second, and more deliberately, 398 of the 1069 frames---about 37 percent---contain
no laser at all. These negatives are not waste; they teach the model what
\emph{absence} looks like, so it stays silent when the dot is off the surface
instead of inventing one. For a controller where a single false detection yanks the
spotlight to the wrong place, suppressing false positives is at least as important
as finding the true dot. We exported the set from Roboflow with a 70/20/10 split
and applied augmentation to the training portion only, which expands it to 2168
training images against 214 validation and 107 test images held out unaugmented.
Training ran on the development GPU with mixed precision. We use the nano model on
purpose: it is small enough to run on the laptop's integrated GPU and even on its
CPU, which is the whole point of the project.

At run time each camera frame goes to the detector; if it returns one or more
boxes we take the highest-confidence one and use its centre as the laser's
camera-pixel position. A modest confidence threshold suppresses spurious
responses, and because we keep only the best box, the detector emits at most one
laser position per processed frame for the renderer to consume asynchronously. We
report accuracy at two input sizes in Section~\ref{sec:results}: the model's
native size, and the smaller size we actually deploy at for speed.

% ---------------------------------------------------------------------------
\subsection{Camera-to-Game Calibration}\label{sec:calib}
% ---------------------------------------------------------------------------
The detector speaks in camera pixels; the game thinks in its own
$1920\times1080$ space. We bridge them with a one-time four-corner calibration.
The operator clicks the four corners of the projected area in the camera image,
in the order top-left, top-right, bottom-right, bottom-left; the matching
game-space corners are the corners of the $1920\times1080$ rectangle. From these
four correspondences we compute a $3\times3$ homography $H$ (with
\texttt{cv2.findHomography}) and map every detected point through it (with
\texttt{cv2.perspectiveTransform}). The homography is saved and reused until the
camera or projector moves. One global homography is enough here because the
surface is flat and the game grid is coarse; the small residual error it leaves is
discussed in Section~\ref{sec:discussion}. A constant pixel offset, adjustable
live with the arrow keys and saved with the calibration, soaks up any systematic
bias between where the laser points and where the spotlight lands.

% ---------------------------------------------------------------------------
\subsection{The Flashlight-Reveal Interaction}\label{sec:reveal}
% ---------------------------------------------------------------------------
The first game's interaction is built on occlusion, not a cursor. Each frame is
drawn normally and then covered by a near-opaque dark overlay; a soft-edged
circular hole is punched in the overlay at the laser position, so only a small
disc around the laser is visible. In practice the overlay is one alpha surface
that we allocate once and reuse, and the spotlight is a pre-built radial texture
subtracted into it at the laser point---details that matter for performance, as
Section~\ref{sec:results} explains. The lit disc has a radius of 160 game pixels
with a soft falloff.

Three small behaviours keep the light from feeling twitchy. It fades in over a
fraction of a second when a detection arrives and fades out when the laser
disappears, so it does not flicker on and off. A detection counts as current only
for a short staleness window; past that the light starts to fade. And an optional
first-order filter on the mapped position trades a little latency for steadiness.
We tuned all of these---radius, fade time, staleness, smoothing, falloff,
darkness---in a stripped-down sandbox that runs the exact camera, detector, and
calibration but draws only the reveal, so the settings carry straight over.

\noindent\textbf{Edge-biased maze.} While testing we noticed the detector was most
reliable near the left and right of the frame and weakest in the centre, a bias we
trace to the training data (Section~\ref{sec:discussion}). Instead of only
gathering more images, we turned it into a level-design rule: the maze the thief
travels is laid out so its main corridors hug the edges, putting the action where
detection is strongest. It is a design response to a model limitation, and it cost
nothing to apply.

% ---------------------------------------------------------------------------
\subsection{Two Thief Controllers}\label{sec:adversaries}
% ---------------------------------------------------------------------------
The thief has to feel like a real opponent: it should head for the treasure,
notice when the light is on it, and try to get away. We built its behaviour two
different ways over the life of the project, and since they make an instructive
contrast we describe both.

\noindent\textbf{Learned: tabular Q-learning.} The first controller treats the
thief as a small Markov decision process. The state is the pair of grid cells
holding the thief and the laser, $(r_t,c_t,r_\ell,c_\ell)$, with $(-1,-1)$ for
``no laser''; the actions are the four grid moves, and a move into a wall is a
penalised no-op. The reward is shaped to want the treasure and fear the light:
reaching the treasure ends the episode with $+1$, being held in the light long
enough ends it with $-1$, a small $-0.01$ per step discourages dawdling, a $-0.05$
discourages bumping walls, and running out of time costs $-0.5$. Because the state
and action spaces are small and discrete, we store the policy as a sparse Q-table
and learn it with ordinary Q-learning ($\alpha=0.1$, $\gamma=0.95$,
$\varepsilon$-greedy decaying from 1.0 to 0.05 over the first 80 percent of
training), which is readable and converges in this finite
case~\cite{watkins1992,sutton2018}. To keep the thief from overfitting one kind of
opponent, each training episode samples one of four scripted laser behaviours: an
idle laser (so it just learns to navigate), a patrolling laser, a human-like laser
that tracks the thief with reaction lag and jitter, and a random laser. We trained
for 20{,}000 episodes on the desktop GPU, which takes only a few seconds and ends
with about 3{,}900 learned states. At play time the table is loaded greedily, with
no further learning. This controller is fast and interpretable, but it learns one
specific maze; change the maze and it has to be retrained.

\noindent\textbf{Search-based: A* with laser avoidance.} The second controller
plans instead of learning. Each round generates a fresh random maze, so a fixed
policy is no use; the thief instead runs A* over the grid with a Manhattan-distance
heuristic~\cite{hart1968} and replans at every cell. The laser enters as a cost,
not as a wall: cells near the current laser position carry an extra traversal
penalty, so the planner naturally bends the route away from the light while still
reaching its goal, and a fallback plan that ignores the laser keeps it from
freezing if the light sits across the only corridor. The game is two-phase. First
the thief heads for the treasure; once it has the treasure the goal switches to
escaping through one of two doors placed far apart and far from the treasure, so a
single spotlight cannot cover both exits. The thief scores the doors as it runs,
adding a penalty when the laser is guarding a door or sitting on the route to it,
and it commits to one only once it has left the shared part of the path---with a
little hysteresis so it does not dither between them. Two reactive touches make it
read as alive: when the light lands on it, it bolts, briefly cutting its
cell-crossing time and quickening its walk animation, and the catch is a hold,
filling while the light stays within 100 pixels of it and draining when it slips
off, so the player must track it rather than tag it once.

\noindent\textbf{What the comparison taught us.} The two controllers trade off
along clear lines. Q-learning gives genuinely emergent evasion---it discovers, for
example, that hugging a wall away from a patrolling laser pays off---but only on
the maze it trained on, and its quality is bounded by how well our four scripted
lasers stand in for a real player. A* needs no training and drops onto any
freshly generated maze, with avoidance that is explicit and easy to reason about,
but its behaviour is only as clever as the costs we hand-write; it does not
\emph{learn} anything about how this particular player aims. The practical reason
we ended up leaning on A* for the random-maze version is mundane but real:
procedurally generated mazes broke the one assumption the trained table depended
on. We think the honest reading is that the learned policy is the better fit for a
fixed, hand-designed level, and the planner is the better fit once levels are
generated on the fly---a point we expand in future work.

% ===========================================================================
\section{The Games}\label{sec:games}
% ===========================================================================

% ---------------------------------------------------------------------------
\subsection{Game One: Thief Hunt}
% ---------------------------------------------------------------------------
The first game puts everything above together. The player sees a black screen and
a small moving pool of light wherever the laser points. Hidden in the dark, a
thief works its way through the maze toward a treasure; the player sweeps the
light to find it and then holds the light on it to catch it. Catching is the dwell
mechanic from Section~\ref{sec:system}: a ring fills while the light stays on the
thief and drains when it comes off, and a full ring is a catch. The thief wins by
reaching the treasure (in the random-maze version, by then escaping through a
door). The treasure itself is hidden in the dark like everything else, which, as a
side benefit, stopped the detector from ever mistaking a permanently lit object
for the laser. The thief is driven by one of the two controllers of
Section~\ref{sec:adversaries}.

% ---------------------------------------------------------------------------
\subsection{Game Two: Fruit Ninja}
% ---------------------------------------------------------------------------
The second game reuses the same laser as a blade. Fruit is thrown up from below
on gravity arcs, and the player slices it by sweeping the laser through it. A
slice is the gesture from Section~\ref{sec:system}: the segment between the
laser's last and current position has to pass within a fruit's radius while the
laser is moving above a speed threshold (we use 22 pixels per frame). The speed
gate is what separates slicing from aiming---resting the laser on a fruit does
nothing, so the player cannot cheat by parking the dot. As the score climbs, fruit
spawns faster and bombs start to appear; slicing a bomb, or letting fruit fall
unsliced, costs a life, and several fruit caught in one fast sweep score a combo.
This game deliberately turns the on-screen pointer marker off, because the laser
dot is itself the visible blade and a second white marker drawn on top of it only
washed out the green the detector was keying on. One laser-specific detail mattered
more than we expected: when the beam blinks out for a frame, we throw away the last
position, otherwise the next detection a moment later would draw a slice straight
across every fruit in between.

The point of describing the two games side by side is how little the second one
cost. Fruit Ninja shares the entire stack below the input layer with Thief
Hunt---same camera, same detector, same calibration, same pointer snapshot---and
adds only its own rules. It needed no change to the vision pipeline at all.

% ---------------------------------------------------------------------------
\subsection{Toward More Games}\label{sec:moregames}
% ---------------------------------------------------------------------------
Two games was never the end goal; it was the proof that the platform carries more
than one. Because every game speaks only to the pointer snapshot and the dwell and
gesture helpers, the genres that fit are the ones a single, occasionally vanishing,
slightly jittery point can drive. Dwell-to-place tower defence is a clean fit: hold
the light on a spot to build. So are whack-a-mole-style sweep games, light-painting
and tracing, simple rhythm games where the laser hits targets on the beat, and
hidden-object or reveal games that reuse the flashlight directly. What does
\emph{not} fit---typing, button-heavy menus, anything needing pixel-perfect
clicks or two cursors---is ruled out by the same input layer that makes the good
cases easy, which is the whole reason we built that layer narrow. Adding one of
these is, in engineering terms, writing a new game module against an interface we
already have, not extending the pipeline.

% ===========================================================================
\section{Implementation and Deployment Challenges}\label{sec:challenges}
% ===========================================================================
A running theme of the project is that the hard parts of a live demo were in
deployment, not in the algorithms. We describe the three that shaped the system
most, each as problem, diagnosis, and fix.

\noindent\textbf{Camera pipeline fragility.} Getting the phone's camera into the
laptop as a virtual webcam was the single biggest time sink. With one popular
phone-as-webcam bridge, OpenCV would variously throw a backend exception, open the
wrong physical device (including the laptop's infrared Windows-Hello camera), or
hand back a single frozen frame. We wrote a small diagnostic that reads many
frames in a row and reports whether the mean brightness \emph{changes} (a live
stream) or stays \emph{constant} (a dead buffer), along with the negotiated pixel
format; it showed the feed was a static buffer with no real stream negotiated,
which we confirmed when an external capture tool reading the same virtual device
also got a frozen frame while the bridge's own window showed live video. We
switched to a phone-as-webcam application that produced a genuine live device, and
we specify a custom push-based camera app as a sturdier replacement in future
work.

\noindent\textbf{Calibration through projected feedback.} When the laptop screen
is mirrored to the projector, any window that shows the camera feed---such as the
calibration view---is itself projected and re-filmed, giving a recursive
``infinity-mirror'' image that makes clicking exact corners hard. We worked around
it by clicking the outermost projected rectangle and by making the tuning and game
tools never display the raw feed on the projected screen. The residual error from
this is the main source of the small spotlight offset in
Section~\ref{sec:discussion}, which the saved constant offset largely cancels.

\noindent\textbf{Making the reveal feel instant.} On the GPU-less laptop the
reveal could lag visibly behind the laser. Why is the subject of
Section~\ref{sec:results}; the short version is that the fix was not one faster
model but a sequence of changes---shrinking the inference input, removing per-frame
surface allocations in the renderer, capping the render and detection rates so the
two threads stop fighting for the CPU, and finally moving inference onto the
integrated GPU.

% ===========================================================================
\section{Experiments and Results}\label{sec:results}
% ===========================================================================
We report three things: how accurately the detector finds the dot, how well the
learned thief plays in simulation, and how the full pipeline behaves in real time
on the demonstration laptop. The last is the paper's main empirical contribution.

\noindent\textbf{Detection accuracy.} Table~\ref{tab:detect} reports the
fine-tuned detector on the held-out test and validation splits at two input sizes.
At the model's native size the test mAP50 is 0.88, with precision and recall both
near 0.85. At the smaller 320-pixel input we actually deploy---chosen because it
roughly halves inference cost on the laptop---test mAP50 falls to about 0.82. That
gap is the accuracy we trade for speed, and on this game it is a trade worth
making: a slightly looser box on a laser dot still gives a good centre, and the
extra frames per second matter far more to how the reveal feels.
Qualitatively, detection was reliable toward the left and right of the frame and
weaker in the centre, which is what motivated the edge-biased maze
(Section~\ref{sec:reveal}).

\begin{table}[t]
\caption{Detector Accuracy (Single-Class Green Laser, Nano YOLO)}
\label{tab:detect}
\begin{center}
\renewcommand{\arraystretch}{1.3}
\begin{tabular}{|p{1.5cm}|p{0.9cm}|c|c|c|c|}
\hline
\textbf{Input} & \textbf{Split} & \textbf{mAP50} & \textbf{mAP50-95} &
\textbf{P} & \textbf{R} \\
\hline
640 (native) & test & 0.881 & 0.429 & 0.851 & 0.847 \\
\hline
640 (native) & val  & 0.951 & 0.483 & 0.919 & 0.931 \\
\hline
\textbf{320 (deployed)} & \textbf{test} & \textbf{0.816} & \textbf{0.395} &
\textbf{0.845} & \textbf{0.831} \\
\hline
320 (deployed) & val  & 0.858 & 0.374 & 0.857 & 0.827 \\
\hline
\end{tabular}
\end{center}
\end{table}

\noindent\textbf{Learned thief in simulation.} Trained for 20{,}000 episodes
against the four scripted laser behaviours, the Q-learning thief reaches the
treasure in 95 percent of simulated rounds and is caught in the other 5 percent.
This is a simulation result against scripted opponents, not a measurement against
human players, and we read it only as evidence that the learned policy navigates
and evades competently on its fixed maze.

\noindent\textbf{Real-time deployment: the inference-versus-render contention.}
Making the reveal feel instant on the GPU-less laptop, not improving accuracy, was
the dominant challenge, and it produced the paper's main finding.
Table~\ref{tab:fps} reports the observed inference rate and the resulting feel for
four configurations at the deployed 320-pixel input.

\begin{table}[t]
\caption{Real-Time Behaviour by Inference Backend on the GPU-less Demonstration
Laptop (Intel Core~i7, Iris Xe iGPU), Inference Input Size 320}
\label{tab:fps}
\begin{center}
\renewcommand{\arraystretch}{1.3}
\begin{tabular}{|p{2.95cm}|p{1.35cm}|p{2.7cm}|}
\hline
\textbf{Configuration} & \textbf{Inference (fps)} & \textbf{Effect on real-time feel} \\
\hline
PyTorch, CPU & $\sim$9--10 & Smoothest of the CPU-only options \\
\hline
PyTorch, CPU, naive full-screen render & $\sim$3 & Rendering and inference contend for the CPU \\
\hline
ONNX Runtime, CPU & $\sim$14--17 & Faster, but uncapped inference starves the renderer (screen $\sim$1 fps) until detection is rate-capped \\
\hline
\textbf{OpenVINO, Iris Xe iGPU} & \textbf{near-instant} & \textbf{Inference offloaded off the CPU; contention removed; smooth} \\
\hline
\end{tabular}
\end{center}
\end{table}

The pattern in Table~\ref{tab:fps} is the result we want to stress. Plain PyTorch
on the CPU was fine on its own, but the moment the full-screen renderer was added
naively---reallocating a large overlay surface and rescaling the background every
frame---inference collapsed to roughly three frames per second. The renderer was
eating the CPU that inference needed. Switching the detector to ONNX Runtime made
inference \emph{faster} in isolation, about 14 to 17 frames per second, and yet the
system felt worse, because the now-faster inference thread ran flat out and starved
the renderer until we capped the detection rate to leave it idle gaps. So on this
machine raw inference speed was not the binding constraint; contention for the one
CPU between inference and rendering was. The configuration that broke the deadlock
was OpenVINO targeting the Intel Iris Xe integrated GPU. By moving inference off
the CPU entirely it left the processor free to render, removed the contention, and
gave the near-instant reveal we were after---the biggest single gain of the
project. Two render-side changes were needed alongside it: reusing one overlay
surface and pre-scaling the background once instead of per frame, and capping the
render and detection rates so neither thread monopolises the CPU. We note that the
performance figures here are observations from run logs on a single laptop, not a
controlled benchmark; we treat them as establishing the qualitative pattern, which
was consistent across sessions, rather than precise constants.

% ===========================================================================
\section{Discussion and Limitations}\label{sec:discussion}
% ===========================================================================
\noindent\textbf{Thermal throttling.} The integrated-GPU configuration was
near-instant, but after roughly a minute of sustained use its responsiveness
drifted back toward the CPU baseline. We put this down to the shared power and
thermal budget of Intel laptop processors: the CPU and the integrated GPU draw from
one budget, so sustained inference heats the package and throttles both. Capping
the detection rate (still instant to the eye, but less heat), running on wall power
in a high-performance profile, and adding active cooling are the mitigations we
identify.

\noindent\textbf{Calibration sensitivity and spotlight offset.} A single global
four-corner homography cannot model lens distortion~\cite{moreno2012}, and the
infinity-mirror calibration (Section~\ref{sec:challenges}) makes exact corner
selection harder. Together they leave a small, position-dependent offset between
the laser and the centre of the revealed spotlight. A constant offset correction
removes most of it, but an offset that grows toward the edges points to
perspective or lens error that only better calibration can fix.

\noindent\textbf{Dataset bias.} The centre-versus-edge asymmetry is a property of
our training data, not of the architecture. We turned it into a level-design rule
with the edge-biased maze, but the principled fix is to retrain on footage shot
under the real demonstration conditions, deliberately including centre-frame and
varied-background cases.

\noindent\textbf{Scope of the controllers and the games.} Our 95 percent figure
for the learned thief is against scripted lasers in simulation, not against people,
and we have not run a formal user study of either game. The A* controller, while it
generalises across mazes, only avoids the light as well as our hand-written costs
tell it to. Both are honest, working controllers; neither has been evaluated
against human players yet, and that evaluation is the obvious next step.

% ===========================================================================
\section{Conclusion and Future Work}\label{sec:conclusion}
% ===========================================================================
We built LaserLearn, a real-time platform that carries a single green laser
pointer from a phone camera, through a fine-tuned single-class detector and a
four-corner homography, into a small laser-faithful input layer that more than one
game can sit on. On that layer we built two games---a flashlight-reveal thief hunt
with two interchangeable adversary controllers, and a laser-slicing arcade
game---and the second cost almost nothing precisely because the platform was
designed for it. Putting the system on a commodity, GPU-less laptop surfaced our
main finding: the limit on real-time feel was not raw inference speed but
contention for the processor between inference and rendering, and offloading
inference to the integrated GPU through OpenVINO removed it and was the largest
single improvement we made.

The clearest next step is more games. The platform was built to host them, and
dwell-to-place tower defence, sweep and reveal games, and simple rhythm games all
map onto the input layer we already have. Beyond that we want to put the two thief
controllers in front of real players and measure them properly; retrain the
detector on footage captured under true demonstration conditions to close the
centre-frame gap; replace the fragile virtual-webcam bridge with a robust
push-based phone-camera app; and explore integer quantization and active cooling to
hold the integrated-GPU speed-up past the onset of thermal throttling.

% ===========================================================================
\begin{thebibliography}{00}

\bibitem{kirstein1998} C. Kirstein and H. M\"uller, ``Interaction with a projection screen using a camera-tracked laser pointer,'' in \emph{Proc. Int. Conf. Multimedia Modeling (MMM)}, Lausanne, Switzerland, 1998, pp. 191--192.

\bibitem{olsen2001} D. R. Olsen Jr. and T. Nielsen, ``Laser pointer interaction,'' in \emph{Proc. SIGCHI Conf. Human Factors in Computing Systems (CHI)}, Seattle, WA, USA, 2001, pp. 17--22.

\bibitem{ahlborn2005} B. A. Ahlborn, D. Thompson, O. Kreylos, B. Hamann, and O. G. Staadt, ``A practical system for laser pointer interaction on large displays,'' in \emph{Proc. ACM Symp. Virtual Reality Software and Technology (VRST)}, Monterey, CA, USA, 2005, pp. 106--109.

\bibitem{lapointe2005} J.-F. Lapointe and G. Godin, ``On-screen laser spot detection for large display interaction,'' in \emph{Proc. IEEE Int. Workshop Haptic Audio Visual Environments and their Applications (HAVE)}, Ottawa, ON, Canada, 2005, pp. 72--76.

\bibitem{krstinic2014} D. Krsti\'ni\'c, A. Kuzmani\'c Skelin, and I. Milati\'c, ``Laser spot tracking based on modified circular Hough transform and motion pattern analysis,'' \emph{Sensors}, vol. 14, no. 11, pp. 20112--20133, 2014.

\bibitem{chavez2012} F. Ch\'avez, F. Fern\'andez, M. J. Gacto, and R. Alcal\'a, ``Hybrid laser pointer detection algorithm based on template matching and fuzzy rule-based systems for domotic control in real home environments,'' \emph{Applied Intelligence}, vol. 36, no. 4, pp. 907--926, 2012.

\bibitem{raskar1998office} R. Raskar, G. Welch, M. Cutts, A. Lake, L. Stesin, and H. Fuchs, ``The office of the future: A unified approach to image-based modeling and spatially immersive displays,'' in \emph{Proc. SIGGRAPH}, Orlando, FL, USA, 1998, pp. 179--188.

\bibitem{raskar1998sar} R. Raskar, G. Welch, and H. Fuchs, ``Spatially augmented reality,'' in \emph{Proc. 1st IEEE Int. Workshop Augmented Reality (IWAR)}, San Francisco, CA, USA, 1998.

\bibitem{raskar2001shaderlamps} R. Raskar, G. Welch, K.-L. Low, and D. Bandyopadhyay, ``Shader lamps: Animating real objects with image-based illumination,'' in \emph{Proc. 12th Eurographics Workshop on Rendering (EGWR)}, London, U.K., 2001, pp. 89--102.

\bibitem{pinhanez2001} C. Pinhanez, ``The everywhere displays projector: A device to create ubiquitous graphical interfaces,'' in \emph{Proc. Int. Conf. Ubiquitous Computing (UbiComp)}, LNCS 2201, Atlanta, GA, USA, 2001, pp. 315--331.

\bibitem{jones2013illumiroom} B. R. Jones, H. Benko, E. Ofek, and A. D. Wilson, ``IllumiRoom: Peripheral projected illusions for interactive experiences,'' in \emph{Proc. SIGCHI Conf. Human Factors in Computing Systems (CHI)}, Paris, France, 2013, pp. 869--878.

\bibitem{jones2014roomalive} B. Jones, R. Sodhi, M. Murdock, R. Mehra, H. Benko, A. Wilson, E. Ofek, B. MacIntyre, N. Raghuvanshi, and L. Shapira, ``RoomAlive: Magical experiences enabled by scalable, adaptive projector-camera units,'' in \emph{Proc. ACM Symp. User Interface Software and Technology (UIST)}, Honolulu, HI, USA, 2014, pp. 637--644.

\bibitem{song2007tangible} P. Song, S. Winkler, and J. Tedjokusumo, ``A tangible game interface using projector-camera systems,'' in \emph{Proc. Int. Conf. Human-Computer Interaction (HCII)}, LNCS 4551, Beijing, China, 2007, pp. 956--965.

\bibitem{moreno2012} D. Moreno and G. Taubin, ``Simple, accurate, and robust projector-camera calibration,'' in \emph{Proc. 2nd Int. Conf. 3D Imaging, Modeling, Processing, Visualization and Transmission (3DIMPVT)}, Zurich, Switzerland, 2012, pp. 464--471.

\bibitem{hartley2004} R. Hartley and A. Zisserman, \emph{Multiple View Geometry in Computer Vision}, 2nd ed. Cambridge, U.K.: Cambridge Univ. Press, 2004.

\bibitem{ridel2014} B. Ridel, P. Reuter, J. Laviole, N. Mellado, N. Couture, and X. Granier, ``The revealing flashlight: Interactive spatial augmented reality for detail exploration of cultural heritage artifacts,'' \emph{ACM Journal on Computing and Cultural Heritage}, vol. 7, no. 2, art. 6, 2014.

\bibitem{bier1993} E. A. Bier, M. C. Stone, K. Pier, W. Buxton, and T. D. DeRose, ``Toolglass and magic lenses: The see-through interface,'' in \emph{Proc. SIGGRAPH}, Anaheim, CA, USA, 1993, pp. 73--80.

\bibitem{watkins1992} C. J. C. H. Watkins and P. Dayan, ``Q-learning,'' \emph{Machine Learning}, vol. 8, no. 3--4, pp. 279--292, 1992.

\bibitem{sutton2018} R. S. Sutton and A. G. Barto, \emph{Reinforcement Learning: An Introduction}, 2nd ed. Cambridge, MA, USA: MIT Press, 2018.

\bibitem{mnih2015} V. Mnih, K. Kavukcuoglu, D. Silver, A. A. Rusu, J. Veness, M. G. Bellemare, et al., ``Human-level control through deep reinforcement learning,'' \emph{Nature}, vol. 518, no. 7540, pp. 529--533, 2015.

\bibitem{tan1993} M. Tan, ``Multi-agent reinforcement learning: Independent vs. cooperative agents,'' in \emph{Proc. 10th Int. Conf. Machine Learning (ICML)}, Amherst, MA, USA, 1993, pp. 330--337.

\bibitem{littman1994} M. L. Littman, ``Markov games as a framework for multi-agent reinforcement learning,'' in \emph{Proc. 11th Int. Conf. Machine Learning (ICML)}, New Brunswick, NJ, USA, 1994, pp. 157--163.

\bibitem{desouza2021} C. de Souza, R. Newbury, A. Cosgun, P. Castillo, B. Vidolov, and D. Kuli\'c, ``Decentralized multi-agent pursuit using deep reinforcement learning,'' \emph{IEEE Robotics and Automation Letters}, vol. 6, no. 3, pp. 4552--4559, 2021.

\bibitem{baker2020} B. Baker, I. Kanitscheider, T. Markov, Y. Wu, G. Powell, B. McGrew, and I. Mordatch, ``Emergent tool use from multi-agent autocurricula,'' in \emph{Proc. Int. Conf. Learning Representations (ICLR)}, 2020.

\bibitem{glavin2015} F. G. Glavin and M. G. Madden, ``Adaptive shooting for bots in first person shooter games using reinforcement learning,'' \emph{IEEE Trans. Computational Intelligence and AI in Games}, vol. 7, no. 2, pp. 180--192, 2015.

\bibitem{hart1968} P. E. Hart, N. J. Nilsson, and B. Raphael, ``A formal basis for the heuristic determination of minimum cost paths,'' \emph{IEEE Trans. Systems Science and Cybernetics}, vol. 4, no. 2, pp. 100--107, 1968.

\bibitem{redmon2016} J. Redmon, S. Divvala, R. Girshick, and A. Farhadi, ``You only look once: Unified, real-time object detection,'' in \emph{Proc. IEEE Conf. Computer Vision and Pattern Recognition (CVPR)}, Las Vegas, NV, USA, 2016, pp. 779--788.

\bibitem{terven2023} J. Terven, D.-M. C\'ordova-Esparza, and J.-A. Romero-Gonz\'alez, ``A comprehensive review of YOLO architectures in computer vision: From YOLOv1 to YOLOv8 and YOLO-NAS,'' \emph{Machine Learning and Knowledge Extraction}, vol. 5, no. 4, pp. 1680--1716, 2023.

\bibitem{gorbachev2019} Y. Gorbachev, M. Fedorov, I. Slavutin, A. Tugarev, M. Fatekhov, and Y. Tarkan, ``OpenVINO deep learning workbench: Comprehensive analysis and tuning of neural networks inference,'' in \emph{Proc. IEEE/CVF Int. Conf. Computer Vision Workshops (ICCVW)}, Seoul, South Korea, 2019.

\bibitem{edgeframeworks2025} ``Comparative analysis of deep learning inference frameworks on edge devices,'' \emph{Electronics}, vol. 14, no. 15, art. 2977, 2025. [Author list to be verified on publisher page.]

\bibitem{edgecontention2022} ``Multi-model running latency optimization in an edge computing paradigm,'' \emph{Sensors}, vol. 22, no. 16, art. 6097, 2022. [Author list to be verified on publisher page.]

\end{thebibliography}

\end{document}
